\documentclass[../main.tex]{subfiles}

\begin{document}

\section{Task 1: Single Burn Arc}\label{sec:task1}

\subsection{Problem Statement}

The simplest mission profile consists of a single continuous powered arc starting from rest at the origin, $(x, y, v_x, v_y) = (0, 0, 0, 0)$ with $m(0) = 1$, and ending at the orbital injection state satisfying Eq.~\eqref{eq:terminal}. The mass-flow rate $Q$ is the primary free parameter swept over the range $Q \in [1.8, 7]$; values below $Q_\mathrm{min} = 1/c \approx 1.67$ would yield insufficient thrust at liftoff ($\left.T/W\right|_{t=0} = cQ < 1$).

\subsection{Boundary Value Problem Formulation}

The four shooting unknowns $\boldsymbol{\zeta}$ (Eq.~\ref{eq:unknowns}) determine the costate initial conditions and the free flight time. The shooting residual is
\begin{equation}
    \mathbf{r}_1(\boldsymbol{\zeta}) =
    \begin{bmatrix}
        y(t_f) - y_f \\[2pt]
        v_x(t_f) - 1 \\[2pt]
        v_y(t_f) \\[2pt]
        \mathcal{H}(0)
    \end{bmatrix} = \mathbf{0}.
    \label{eq:residual1}
\end{equation}
The first three components enforce the terminal state constraints \eqref{eq:terminal}; the fourth is the free-time condition $\mathcal{H} = 0$, imposed at $t_0$ through the algebraic form \eqref{eq:H0}. The mass-costate transversality $\lambda_m(t_f) = 1$ and the terminal value of $\mathcal{H}$ are no longer needed as residuals, having been folded into the $\lambda_{m,0} = 1$ normalization and the autonomy of the system.

\subsection{Velocity-Loss Decomposition}

The total velocity loss $W_\mathrm{tot} = \Delta V_\mathrm{ideal} - v_f$ compares the ideal Tsiolkovsky delta-V,
\begin{equation}
    \Delta V_\mathrm{ideal} = c\ln\!\frac{m_0}{m_f} = c\ln\!\frac{1}{m_f},
\end{equation}
against the achieved injection speed $v_f = 1$. It is decomposed as
\begin{align}
    W_d &= \int_0^{t_f} \frac{T}{m}\bigl[1 - \cos(\varphi - \psi)\bigr]\,\mathrm{d}t
         && \text{(thrust misalignment loss)},
    \label{eq:Wd}\\
    W_g &= \int_0^{t_f} \sin\psi\,\mathrm{d}t
         && \text{(gravity loss)},
    \label{eq:Wg}
\end{align}
where $\psi = \arctan(v_y/v_x)$ is the flight-path angle. Misalignment loss $W_d$ arises from the fraction of thrust not aligned with the velocity vector; gravity loss $W_g$ is the work done against gravity projected along the trajectory.

\subsection{Results}\label{sec:task1_results}

Figure~\ref{fig:task1a} shows the final mass as a function of $Q$ for the three target altitudes. Each curve exhibits a clear interior maximum $Q^*$: at low $Q$, the long burn time accumulates large gravity losses; at high $Q$, the vehicle climbs steeply and thrust-misalignment loss dominates.

\begin{figure}[H]
    \centering
    \includegraphics[width=0.78\textwidth]{task1_task_1a_final_mass_vs_q}
    \caption{Task 1a: final mass $m_f$ as a function of the mass-flow rate $Q$ for three target altitudes $y_f$.}
    \label{fig:task1a}
\end{figure}

The interior maxima are quantified in Table~\ref{tab:qstar}: the optimal rate $Q^\ast$ decreases from $2.52$ to $2.19$ as the target altitude rises from $y_f = 0.04$ to $0.06$, and the delivered payload turns \emph{negative} at $y_f = 0.06$---a single stage cannot reach that orbit with positive payload, which is precisely what motivates the staging of Task~4. The nominal $Q = 2$ of the assignment lies within a few percent of $Q^\ast$.

\begin{table}[H]
\centering
\caption{Payload-optimal mass-flow rate $Q^\ast$ (maximizing the final mass) for the three target altitudes, read from the $Q$-sweep of Figure~\ref{fig:task1a}. The delivered single-stage payload is $m_\mathrm{pay} = m_f(1+\eta) - \eta$.}
\label{tab:qstar}
\begin{tabular}{cccc}
\toprule
$y_f$ & $Q^\ast$ & $m_f^\ast$ & $m_\mathrm{pay}^\ast$ \\
\midrule
$0.04$ & $2.52$ & $0.1180$ & $+0.0298$ \\
$0.05$ & $2.33$ & $0.1012$ & $+0.0114$ \\
$0.06$ & $2.19$ & $0.0854$ & $-0.0060$ \\
\bottomrule
\end{tabular}
\end{table}

Figure~\ref{fig:task1b} confirms this trade-off by decomposing the velocity loss at $y_f = 0.04$. The gravity loss $W_g$ is large at low $Q$ and decreases monotonically with $Q$, while the misalignment loss $W_d$ has a broad minimum near $Q^*$.

\begin{figure}[H]
    \centering
    \includegraphics[width=0.78\textwidth]{task1_task_1b_velocity_losses}
    \caption{Task 1b: velocity-loss decomposition---misalignment $W_d$, gravity $W_g$, and total $W_\mathrm{tot}$---at $y_f = 0.04$.}
    \label{fig:task1b}
\end{figure}

Figures~\ref{fig:task1c_traj} and~\ref{fig:task1c_angles} show the optimal trajectory and control history at $Q = Q^*$ for $y_f = 0.04$. The thrust angle $\varphi$ decreases smoothly from near-vertical at launch to near-horizontal at injection, following the bilinear-tangent law \eqref{eq:bilinear}. The velocity angle $\psi$ tracks $\varphi$ closely once the initial transient has died out.

\begin{figure}[H]
    \centering
    \begin{subfigure}[b]{0.48\textwidth}
        \centering
        \includegraphics[width=\textwidth]{task1_task_1c_trajectory}
        \caption{Optimal trajectory in the $(x,y)$ plane.}
        \label{fig:task1c_traj}
    \end{subfigure}
    \hfill
    \begin{subfigure}[b]{0.48\textwidth}
        \centering
        \includegraphics[width=\textwidth]{task1_task_1c_thrust_and_velocity_angles}
        \caption{Thrust angle $\varphi$ and velocity angle $\psi$ vs.\ time.}
        \label{fig:task1c_angles}
    \end{subfigure}
    \caption{Task 1c: optimal trajectory and control at $Q = Q^*$, $y_f = 0.04$.}
    \label{fig:task1c}
\end{figure}

\end{document}
